\section{Equivalences and reductions}\label{sec:reductions}

This section summarizes archived reductions. The research-record entries use definitions and hypotheses supplied in their linked claim-index evidence; they are not standalone theorem statements or proofs in this manuscript. They are not used to prove the main results. Some restate parts of the open problem in other terms. An equivalence earns no reduction in difficulty by itself. We record these results because they fix the exact shape of the remaining quantifiers.

\begin{researchrecord}[\cid{B.9}; written proof from published and preprint sources; an equivalence]\label{thm:o1}
Let $S_0=\{2,3,5,2243,4547,5843\}$. For every good ordinary prime $p\notin S_0$,
\[
h_p=0\iff\Sha[p^\infty]=0 ,
\]
where $h_p$ is the witness depth of Ota's theta elements \cite{ota} over admissible pairs of auxiliary primes.
\end{researchrecord}
The forward direction is Ota's Theorem~6.4. The converse combines three inputs:
\begin{itemize}[leftmargin=*]
\item Castella--Sano \cite{castellasano}, whose Theorem~B produces a $p$-unit Kurihara number;
\item C.-H.~Kim \cite{kimch}, Theorem~1.10, which localizes the Selmer group at the minimal pair;
\item a normalization check between the period conventions.
\end{itemize}
Uniform primitive witnesses on the ordinary tail are exactly uniform vanishing there. The theorem gives no converse at positive depth or at supersingular primes.

\begin{researchrecord}[\cid{B.10}; written proof; an equivalence]\label{thm:divfield}\leavevmode\par
Let $p\ge7$ with $p\notin\{19,59,487,2243,4547,5843\}$, and put
\[
\mathcal C_p=\mathrm{Hom}_{\GL_2(\F_p)}\bigl(\mathrm{Cl}(\Q(E[p]))/p,\ E[p]\bigr).
\]
Then $\Sha[p]=0$ if and only if $\dim_{\F_p}\mathcal C_p\le2$. Moreover $\dim\mathcal C_p\ge1$.
\end{researchrecord}
The excluded primes limit where this simplified reduction applies. They are \emph{not} a bound on the primes dividing $\#\Sha$.

\begin{researchrecord}[\cid{B.11}; written proof, finite algebra Lean-checked; conditional]\label{thm:line}
\emph{LINE} is the hypothesis that one rational scalar $a/b$ relates the tame leading terms $\delta_n$ and the regulator terms $\rho_n$ at all required auxiliary pairs and primes, by $b\delta_n=a\rho_n$. Under LINE, the primary components at the covered good primes are finite, and they vanish away from the support of the numerator.

\emph{Witness content} weakens LINE to the existence of a nonzero integer $A\in J_E$, and the same conclusion holds under it.
\end{researchrecord}
The finite algebra of both implications is formalized in Lean (the frozen archives A1 and A2). The arithmetic antecedents are not proved: neither the cross-prime scalar of LINE nor the nonzero integer $A$ is known to exist.

LINE is related to the residue observation of Remark~\ref{rem:singletons}, but it is not equivalent to it. LINE needs a common rational comparison on core pairs at every $p$-power level, with auxiliaries that carry $p$-power point torsion; the observation concerns nonvanishing modulo $p$ at selected auxiliaries (Research record~\ref{thm:g2}). Neither is known to imply the other. At the grade of a conjecture (Conjecture~2.3 of Burns--Kurihara--Sano \cite{bks}, in the production conventions) the constant of the tame comparison is
\[
K_E=\frac{4\,c_\infty\,\#\Sha\,\mathrm{Tam}}{\#E(\Q)_{\mathrm{tors}}^2},\qquad c_\infty=\#\pi_0(E(\mathbf{R})).
\]
It matches exactly on six rank-two curves with trivial $\Sha$: $68$ of $68$ rows with nonzero tame regulator, and $7$ of $7$ rows where both sides vanish. The factor $\#\Sha$ was tested on the curve 194040co6, whose analytic order of $\Sha$ is $4$: $17$ of $17$ rows match, and its isogenous curves with trivial $\Sha$ match $40$ of $40$. Every constant there is a power of $2$. No rank-two curve of conductor below $500{,}000$ has odd analytic order of $\Sha$ other than $1$, so an odd factor such as the $9$ expected here cannot be tested on such curves. For $E$ one has $K_E=32\,\#\Sha$, and $288=32\cdot9$ is consistency evidence, at conjecture grade, for $\#\Sha\equiv9$ modulo each observed prime. It says nothing about $q$.

\begin{proposition}[\cid{B.12}; standard deduction with an exact computation]\label{prop:localsel}
\begin{enumerate}[label=(\roman*)]
\item For $p\in\{2,3,5,7,11\}$, the rational $p$-adic Selmer group is\[\Qp\delta_p(P_1)\oplus\Qp\delta_p(P_2).\]
\item For $p=5,7,11$, one has $\Sel_{p^n}(E/\Q)\simeq(\Z/p^n)^2$ for all $n$, and
\[
\ker(\mathrm{loc}_p)=\Qp\bigl(\log_p(P_2)\,\delta_p(P_1)-\log_p(P_1)\,\delta_p(P_2)\bigr).
\]
\item At $p=11$, with $S=\{\infty,2,3,11,2243,4547,5843\}$, the kernel of $\mathrm{loc}_S$ on $\Sel_{11}(E/\Q)$ is the line $\F_{11}\,\delta_{11}(2P_1-P_2)$.
\end{enumerate}
\end{proposition}

\begin{researchrecord}[\cid{B.13}; written source composition; conditional]\label{thm:mcs}
Let $p\ge5$ be a good ordinary prime with $p\notin\{19,59,487\}$ and $p\nmid\#\widetilde E(\F_p)$. Combining Burungale--Castella--Kim \cite{bck}, Kataoka--Sano and Macias Castillo--Sano \cite{mcs} gives two things over a suitable imaginary quadratic field:
\begin{itemize}[leftmargin=*]
\item a base-layer Heegner Stark-system basis;
\item Fitting-ideal identities for the classical Selmer group at every depth.
\end{itemize}
For $K=\Q(\sqrt{-23})$ (Theorem~\ref{thm:d23}) the identity reads
\[
\mathrm{Fitt}^{\,i}\bigl(\Sha(E/K)[p^n]^*\bigr)=I_{i+3}(\epsilon^{\mathrm{Hg}}_{n,0}),
\]
at $p=13,31,41,47,71,73,101,127,139,151$ among the primes $p\le160$.
\end{researchrecord}
The ideals $I_{i+3}$ are not evaluated at any prime, and the Heegner $\Lambda$-index does not determine them (an explicit counterexample is recorded). Moreover, $I_3=\F_{13}$ would require $\Sha(E^{-23}/\Q)[13]=0$ as well as the seed's $13$-part. The twist part is not supplied, so no unit specialization is known at $13$ (see the withdrawn claim \cid{W.1}).

\begin{lemma}[\cid{B.14}; Lean-checked]\label{lem:minors}
Let $A$ be a fixed integer matrix. Suppose that for every $p$, $I_3(p)=\F_p$ holds if and only if some $k\times k$ minor of $A$ is nonzero modulo $p$, and that this holds for at least one $p$. Then $C_1=|\gcd(\text{$k\times k$ minors of }A)|$ is nonzero, and $I_3(p)\ne\F_p$ implies $p\mid C_1$.
\end{lemma}
No such integral model has been constructed, so Lemma~\ref{lem:minors} currently has no instance.

\begin{researchrecord}[\cid{B.16}; written proof applying published constructions]\label{thm:g2}
Let $S_0=\{2,3,5,2243,4547,5843\}$. For every prime $p\notin S_0$ there are infinitely many primes $\ell\nmid pN$ with
\[
\ell\equiv1+p\pmod{p^2},\qquad \bar\rho_{E,p}(\mathrm{Frob}_\ell)=-1,\qquad R_{p,\ell}\ne0\ \text{in }\F_p ,
\]
where $R_{p,\ell}$ is the normalized determinant of the tame pairing on the basis of Theorem~\ref{thm:mw}. In particular $v_p(\ell-1)=1$ and $\#E(\F_\ell)\equiv4\pmod p$.
\end{researchrecord}
The proof fixes one Poincar\'e biextension over $\Q$ before $p$ or $\ell$ is chosen, following Tan \cite{tan} and Jossen \cite{jossen}, shows that the relevant Kummer image is full, and applies Chebotarev's theorem. It is existential and gives no bound for the smallest $\ell$. It does not identify $C_2(p,\ell)$ with $R_{p,\ell}$, so it proves nothing about $\Sha$.

\begin{researchrecord}[\cid{B.17}; written proof from preprint sources]\label{thm:katofact}
Let $p$ be a prime outside $\{2,3,5,19,59,487,2243,4547,5843\}$, let $m$ be a squarefree product of primes $\ell\nmid pN$ with $\ell\equiv1\pmod p$ and $p\nmid\#E(\F_\ell)$, and let $I$ be the augmentation ideal of $\Z_p[\Gal(F/\Q)]$ for the corresponding abelian $p$-extension $F$. Put $s_p=2+\operatorname{corank}\Sha[p^\infty]$ and $A_p=H^2_f(\Q,T_pE)_{\mathrm{tors}}$. If $s_p>2$, the degree-two witness $t_{p,m}$ is $0$. If $s_p=2$, then $\Sha[p^\infty]$ is finite, $\#A_p=\#\Sha[p^\infty]$, and
\[
t_{p,m}=u_{p,m}\,x_{\Q,\Sigma}\,\#A_p\,\det(B_m)\quad\text{in }I^2/I^3,
\]
where $u_{p,m}$ is a unit, $x_{\Q,\Sigma}$ is the coordinate of the determinantal Kato element and $\det(B_m)$ is the Bockstein determinant.
\end{researchrecord}
The inputs are the determinant constructions of Burns--Kurihara--Sano \cite{bks}, arithmetic Selmer duality and a Bockstein theorem, all imported. This is an identity in a finite quotient: a sufficiently divisible coefficient is $0$, and its depth is not a finite truncation endpoint. Nothing here bounds the multiplier $x_{\Q,\Sigma}\,\#A_p\,\det(B_m)$ uniformly in $p$. A nonzero singleton witness in this range forces both $\Sha[p^\infty]=0$ and primitivity of the Kato element; the range excludes $19$.

\begin{researchrecord}[\cid{B.19}; written proof from published and preprint sources]\label{thm:w}
Let $p\notin B^*=\{2,3,19,59,487,2243,4547,5843\}$, and let $\ell\neq\ell'$ be primes with $\gcd(\ell\ell',pN)=1$, $\ell\equiv\ell'\equiv1\pmod p$, $p\nmid\#E(\F_\ell)\,\#E(\F_{\ell'})$ and $R_{p,\ell}R_{p,\ell'}\not\equiv0\pmod p$. With one period, one saturated basis and matching generators,
\[
\frac{C_2(p,\ell)}{(a_\ell-2)R_{p,\ell}}=\frac{C_2(p,\ell')}{(a_{\ell'}-2)R_{p,\ell'}}\quad\text{in }\F_p ,
\]
and both sides vanish if $\Sha[p]\ne0$.
\end{researchrecord}
This is a statement at one prime $p$. It relates no two primes, does not show that the common value is nonzero, and does not derive the value $288$. At $p=5$ it only compares eligible pairs, and it says nothing at $19$, $59$ or $487$. A mismatched generator changes the ratio by a square factor, which can be $-1$; a change of period scales every ratio by one common constant.

\subsection{Kurihara witnesses at every good prime}\label{subsec:kurihara}
Fix a good prime $p\ge5$. Following Castella--Sano \cite[\S1.1]{castellasano} and C.-H.~Kim \cite{kimch}, let $n$ run over the squarefree products of primes $\ell\nmid N$ with $\ell\equiv1\pmod p$ and $a_\ell\equiv\ell+1\pmod p$, and let $\nu(n)$ be the number of prime factors of $n$. The Kurihara number $\delta_n$ is a sum of modular symbols normalized by $\Omega^+_E=\int_{E(\mathbf R)}\omega_E=2\omega_1$ and weighted by discrete logarithms; it lies in a quotient $\Z/I_n$ with $I_n\subseteq p\Z$ for $n\ne1$. Write $M_\infty(\delta)$ for the limit of the minimal $p$-adic orders of the $\delta_n$ \cite[\S1.1.3]{castellasano}. These auxiliary primes are not those of Research record~\ref{thm:g2}.

\begin{researchrecord}[\cid{B.20}; written application of preprint theorems; derived by Codex (OpenAI) and independently reviewed by ChatGPT Pro (OpenAI), with source and deduction checks by Claude models (Anthropic); not a human review; not formalized]\label{thm:kurexist}
For every prime $p\ge5$ of good reduction for $E$, including the supersingular primes and $19$, $59$ and $487$, one has $M_\infty(\delta)=0$. Hence some squarefree Kato conductor $n$, which may depend on $p$, gives a mod-$p$ unit Kurihara witness: $\delta_n\not\equiv0\pmod p$. The bad primes $2243$, $4547$ and $5843$ and the primes $p\le3$ are not covered.
\end{researchrecord}
The supersingular primes use Theorem~A of Castella--Liu--Wan \cite{clw} in its author-hosted revision dated 21 September 2026. Its hypothesis (ram), a prime $\ell\parallel N$ at which the residual representation is ramified, holds with $\ell=2243$, where $E$ has split multiplicative reduction of type $\mathrm I_1$. Moreover $k=2<p$, and $p$-regularity holds because $a_p=0$. At $k=2$ the component proved there is the cyclotomic one for $T_pE$ \cite[Remark~2.20]{clw}. Version~1 of \cite{clw} on arXiv required squarefree level; the revision does not. With Theorem~A and Corollary~A of Castella--Sano \cite{castellasano}, this gives $M_\infty(\delta)=0$, since $\mathrm{Tam}=16$ and the Manin constant, which divides $2^9$ (Theorem~\ref{thm:manin}), are prime to $p$.

The ordinary primes are covered by Theorem~B of Castella--Sano as printed. Its hypotheses are $p>3$, (sur) and $p\nmid c_\phi$, and its proof cites the ordinary main conjecture ([Wan15, BCS25] in \cite{castellasano}). The primes $19$, $59$ and $487$ are ordinary ($a_{19}=4$, $a_{59}=-4$, $a_{487}=-28$), and Theorem~B has no hypothesis that excludes them. The Euler factors at $2243$ and $5843$ are nonunits at these primes ($2242=2\cdot19\cdot59$, $5844=12\cdot487$); both formulations retain these factors and neither cancels them, so they create no extra exception. The supersingular case of Theorem~B requires squarefree level, which is why the supersingular primes use the revision of \cite{clw}.

The revision of \cite{clw} is author-hosted and was not on arXiv on 30 September 2026: arXiv 2608.29470 then listed only version~1, the squarefree-level version. The revision used here is pinned by its SHA-256 hash, recorded in the claim index; a later version could change the hypotheses of Theorem~A. Research record~\ref{thm:kurexist} bounds neither the length of the witness nor anything uniformly in $p$.

\begin{researchrecord}[\cid{B.21}; written deduction from preprint theorems; derived by Codex (OpenAI) and independently reviewed by ChatGPT Pro (OpenAI), with source and deduction checks by Claude models (Anthropic); not a human review; not formalized]\label{thm:rp}
For a good prime $p\ge5$ let $r_p=\min\{\nu(n):\delta_n\not\equiv0\pmod p\}$ be the minimal length of a mod-$p$ unit Kurihara witness; it is finite by Research record~\ref{thm:kurexist}. Then
\[
r_p=\dim_{\F_p}\Sel_p(E/\Q),
\]
$r_p$ is even and at least $2$, and $r_p=2$ if and only if $\Sha(E/\Q)[p^\infty]=0$. This also excludes a hypothetical divisible part; no finiteness of $\Sha$ is assumed.
\end{researchrecord}
The proof uses only the clauses of Theorem~1.8 of \cite{kimch} that do not assume $\Sha$ finite (the corank and the higher Fitting ideals), the parity forced by the root number $+1$, and Nakayama's lemma for the dual Kummer sequence $0\to\Sha[p^\infty]^\vee\to X_p\to\Z_p^2\to0$, which splits; here $X_p$ is the Pontryagin dual of $\Sel_{p^\infty}(E/\Q)$. The index $r_p$ is the first index of a unit, not the first index of a nonzero $p$-adic value. This is an identity, not an upper bound. The abstract dual module $X=\Z_p^2\oplus(\Z/p)^2$ needs four generators, with $\mathrm{Fitt}_2(X)=(p^2)$, $\mathrm{Fitt}_3(X)=(p)$ and $\mathrm{Fitt}_4(X)=\Z_p$. It is compatible with $M_\infty=0$ and first unit index four. This is a module countermodel to an inference from primitivity alone, not a claim about the seed at an uncomputed prime.

\begin{researchrecord}[\cid{B.22}; written deduction from a preprint theorem; conditional; derived by Codex (OpenAI) and independently reviewed by ChatGPT Pro (OpenAI), with source and deduction checks by Claude models (Anthropic); not a human review; not formalized]\label{thm:depth}
Let $p\ge5$ be a good prime, let $K/\Q$ be a real cyclic extension of degree $p$ whose actual conductor $m$ is prime to $pN$ ($m$ may be composite), and let $\chi$ be a faithful character of $\Gal(K/\Q)$. Let $\theta_K^{\mathrm{MT}}$ be the full Mazur--Tate element of $K$ at its actual conductor $m$, normalized as in \cite[Def.~2.2]{bh} with the total real period $\Omega^+=2\omega_1$. If
\[
v_{\zeta_p-1}\bigl(\chi(\theta_K^{\mathrm{MT}})\bigr)<4,
\]
then $\Sha(E/\Q)[p^\infty]=0$.
\end{researchrecord}
The inputs are Hypothesis~1.2 and Theorem~1.4(b) of \cite{bh}, which place $\theta_K^{\mathrm{MT},\#}$ in the zeroth Fitting ideal of the dual Selmer group over $K$ without assuming $\Sha$ finite; Research record~\ref{thm:rp}; and restriction to $\Q$, which is injective because $E(K)[p]=0$ (the image of the Galois group of $K$ contains $\mathrm{SL}_2(\F_p)$). The inequality is strict. A bound $v\le4$ is not enough: over $R=\Z_p[C_p]$ with $t=\gamma-1$, the module $(R/(t))^2\oplus(R/(p,t))^2$ has $t^4$ in its Fitting ideal while its coinvariants need four generators. The element is taken at the actual conductor of $K$; an element at an enlarged conductor must be related to it by the norm relations of \cite[Prop.~2.4]{bh}, with their Euler factors retained. Nonvanishing of a complex $L$-value is not a depth bound. The criterion is sufficient only: no converse from $\Sha(E/\Q)[p^\infty]=0$ to the existence of such a character value is claimed here. The converse holds by \cid{B.23} (derived 2026-09-30); see Research record~\ref{thm:cmp}. It has not been evaluated at any $p\ge53$.

\subsection{Depth, witness length and Selmer dimension}\label{subsec:cmp}
For a good prime $p\ge5$ put $d_p=\dim_{\F_p}\Sel_p(E/\Q)$, let $r_p$ be as in Research record~\ref{thm:rp}, and let $h_p$ be the minimum of $v_{\zeta_p-1}(\chi(\theta_K^{\mathrm{MT}}))$ over the real cyclic extensions $K/\Q$ of degree $p$ at their actual conductors prime to $pN$ and over the faithful characters $\chi$ of $\Gal(K/\Q)$, where $\theta_K^{\mathrm{MT}}$ is the full Bullach--Honnor element; $h_p=+\infty$ is allowed.

\begin{researchrecord}[\cid{B.23}; written deduction from a preprint theorem; derived by ChatGPT Pro (OpenAI) and independently checked by Codex (OpenAI), with a cross-family check by Claude (Anthropic); not a human review; not formalized]\label{thm:cmp}
For every good prime $p\ge5$,
\[
\min(h_p,p-1)=\min(r_p,p-1)=\min(d_p,p-1).
\]
Consequently, for every good prime $p\ge5$,
\[
h_p<4\iff h_p=2\iff r_p=2\iff d_p\le3\iff\Sha(E/\Q)[p^\infty]=0 .
\]
\end{researchrecord}
This is an equivalence between existence statements. It does not say that a preselected field, or every field, attains depth two. It supplies no witness, no new vanishing prime and no upper bound. The strict inequality $<4$ and the cap $p-1$ are part of the statement. The proof applies Theorem~1.4(b) of \cite{bh} over the real abelian $p$-extensions $L$ of conductor prime to $pN$, with Fitting base change and a minimal presentation over the local ring $\F_p[\Gal(L/\Q)]$: since $E(L)[p]=0$, this gives $\theta_L\bmod p\in J^{d_p}$ for the augmentation ideal $J$. The conductor norm relation \cite[Prop.~2.4]{bh}, whose factor at a Kato prime is $-(u-1)^2/u\in J^2$, then identifies the degree-$d_p$ initial form at a minimum unit witness as $\delta_n\prod T_i$. Coefficients divisible by $p$ have valuation at least $p-1$, which is why the equality is capped. Research record~\ref{thm:depth} is the direction $h_p<4\Rightarrow\Sha(E/\Q)[p^\infty]=0$; the converse comes from the following bridge.

\begin{researchrecord}[\cid{B.24}; written deduction from a preprint theorem; conditional; derived by ChatGPT Pro (OpenAI) and independently checked by Codex (OpenAI), with a cross-family check by Claude (Anthropic); not a human review; not formalized]\label{thm:bridge}
Let $p\ge5$ be a good prime, let $q_1\ne q_2$ be Kato primes for $p$, and suppose that $\delta_{q_1q_2}$ is a mod-$p$ unit. Let $L$ be the $C_p^2$ compositum of the degree-$p$ subfields of $\Q(\mu_{q_1})$ and $\Q(\mu_{q_2})$, and let $J=(T_1,T_2)$ be the augmentation ideal of $\F_p[\Gal(L/\Q)]$. In the total-real-period normalization,
\[
\theta_L\equiv\delta_{q_1q_2}\,T_1T_2\pmod{p,\ J^3}.
\]
Every cyclic degree-$p$ quotient of $L$ that is nontrivial on both inertia groups (the diagonal quotient is one) has actual conductor $q_1q_2$, and each of its faithful characters has exact depth $2$.
\end{researchrecord}
The antecedent is a sufficient witness, not an established fact; no two-prime unit witness is supplied here. Here $\delta$ uses the symbol of Castella--Sano normalized by $\Omega_E^+=2\omega_1$, which is one half of the production symbol for this seed.
